\exercisesection{Extension Modules}

¶ 1) An A-complex C is said to be pure if for every right A-module P one has H(P ⊗\_A C) = 0.

\begin{enumerate}
\def\labelenumi{\alph{enumi})}
\item
  A complex homotopic to zero is pure; a complex of flat modules, with zero homology and bounded on the right, is pure. Give an example of a complex of free modules of rank 1, with zero homology but not pure.
\item
  Show that the following conditions are equivalent:
\end{enumerate}

α) C is pure.

β) For every finitely presented A-module M, one has H(Homgr\_A (M, C)) = 0.

\(\gamma)\) There exists a filtered inductive system \((C^{(\alpha)})_{\alpha \in j}\) of complexes homotopic to zero such that \(C = \varinjlim C^{(\alpha)}\) .

(To see that \(\beta\) ) implies \(\gamma\) ), reduce to the case where \(C_i = 0\) for \(i < i_0\) and write \(C_{i_0}\) as an inductive limit of finitely presented modules.)

\begin{enumerate}
\def\labelenumi{\arabic{enumi})}
\setcounter{enumi}{1}
\tightlist
\item
  Let M be an A-module. We say that a submodule \(\mathbf{M}'\) of M is pure if the complex defined by the exact sequence \(0 \to \mathbf{M}' \to \mathbf{M} \to \mathbf{M} / \mathbf{M}' \to 0\) is pure (Exercise 1).
\end{enumerate}

\begin{enumerate}
\def\labelenumi{\alph{enumi})}
\item
  Show that when A is a principal ideal ring, the notion of pure submodule coincides with that of VII, § 2, Exercise 7.
\item
  Suppose that the A-module M is flat. Show that the following conditions are equivalent: \(\alpha\) the submodule M' is pure;
\end{enumerate}

β) the quotient M/M' is flat;

\(\gamma)\) for every ideal \(\mathfrak{a}\) of \(\mathbf{A}\) , we have \(\mathbf{M}' \cap \mathfrak{a}\mathbf{M} = \mathfrak{a}\mathbf{M}'\) .

\begin{enumerate}
\def\labelenumi{\arabic{enumi})}
\setcounter{enumi}{2}
\tightlist
\item
  Let M be an A-module, n an integer ≥ 1. Show that the following conditions are equivalent:
\end{enumerate}

α) M admits a finite n-presentation (X, p.~180, Exercise 6).

β) For every filtered inductive system \((\mathrm{N}_{\alpha})_{\alpha\in I}\) , the canonical homomorphism

\[
\varinjlim \operatorname{Ext} _ {\mathbf {A}} ^ {i} (\mathbf {M}, \mathbf {N} _ {\alpha}) \rightarrow \operatorname{Ext} _ {\mathbf {A}} ^ {i} (\mathbf {M}, \varinjlim \mathbf {N} _ {\alpha})
\]

is bijective for i \textless{} n.

γ) For every family \((\mathbf{P}_{\beta})_{\beta \in J}\) of right A-modules, the canonical homomorphism

\[
\operatorname{Tor} _ {i} ^ {\mathbf {A}} \left(\prod_ {\beta \in J} \mathrm{P} _ {\beta}, \mathrm{M}\right)\rightarrow \prod_ {\beta \in J} \operatorname{Tor} _ {i} ^ {\mathbf {A}} \left(\mathrm{P} _ {\beta}, \mathrm{M}\right) \quad \text { is   bijective   for } \quad i <   n.
\]

δ) For every set I, the canonical homomorphism \(A_{d}^{I} \otimes_{A} M \to M^{I}\) is bijective, and we have

\[
\operatorname{Tor} _ {i} ^ {\mathbf {A}} \left(\mathrm{A} _ {d} ^ {\mathrm{I}}, \mathrm{M}\right) = 0 \quad \text { for } \quad 0 <   i <   n.
\]

(Argue by induction on n, using Exercise 18, p.~170.)

\begin{enumerate}
\def\labelenumi{\arabic{enumi})}
\setcounter{enumi}{3}
\tightlist
\item
  Let P be a complex of projective A-modules, zero on the right, and n an integer such that \(H_{i}(P)=0\) for i\textless n.~Let M be a left A-module, N a right A-module. Show that we have \(H^{i}(\mathrm{Homgr}_{\mathbf{A}}(\mathbf{P},\mathbf{M}))=0\) and \(H_{i}(\mathbf{N}\otimes_{\mathbf{A}}\mathbf{P})=0\) for i\textless n, and that the canonical homomorphisms
\end{enumerate}

\[
\lambda^ {n}: \mathrm{H} ^ {n} (\operatorname{Homgr} _ {\mathbf {A}} (\mathrm{P}, \mathrm{M})) \rightarrow \operatorname{Hom} _ {\mathbf {A}} \left(\mathrm{H} _ {n} (\mathrm{P}), \mathrm{M}\right), \quad \gamma^ {n}: \mathrm{N} \otimes_ {\mathbf {A}} \mathrm{H} _ {n} (\mathrm{P}) \rightarrow \mathrm{H} _ {n} (\mathrm{N} \otimes_ {\mathbf {A}} \mathrm{P})
\]

are bijective.

¶ 5) Let P be a projective A-complex, zero on the right, n an integer ≥ 0. Show that the following conditions are equivalent:

α) There exists a complex \(\mathbf{P}'\) of finitely generated projective A-modules, such that \(\mathbf{P}_i' = 0\) for \(i > n\) or \(i < 0\) , and a morphism \(f: \mathbf{P}' \to \mathbf{P}\) such that \(\mathbf{H}_i(f)\) is bijective for \(i < n\) and surjective for \(i = n\) .

β) For every filtered inductive system of A-modules \((M_{\alpha})_{\alpha \in I}\) , the canonical homomorphism

\[
\varinjlim \mathrm{H} ^ {i} (\operatorname{Homgr} _ {\mathrm{A}} (\mathrm{P}, \mathrm{M} _ {\alpha})) \rightarrow \mathrm{H} ^ {i} (\operatorname{Homgr} _ {\mathrm{A}} (\mathrm{P}, \varinjlim \mathrm{M} _ {\alpha}))
\]

is bijective for \(i < n\) and injective for \(i = n\) .

γ) For any family of right A-modules \((\mathbf{N}_{\alpha})_{\alpha\in I}\) , the canonical homomorphism

\[
\mathrm{H} _ {i} ((\prod_ {\alpha} \mathrm{N} _ {\alpha}) \otimes_ {\mathrm{A}} \mathrm{P}) \rightarrow \prod_ {\alpha} \mathrm{H} _ {i} (\mathrm{N} _ {\alpha} \otimes \mathrm{P})
\]

is bijective for \(i < n\) and surjective for \(i = n\) .

(To prove that \(\alpha\) ) implies \(\beta\) ) and \(\gamma\) ), apply Exercise 4 to the cone of \(f\). To prove that \(\beta\) )

or γ) implies α), reason by induction on n, applying Exercise 4, Exercise 18, p.~170 and Exercise 10, p.~174.)

\begin{enumerate}
\def\labelenumi{\arabic{enumi})}
\setcounter{enumi}{5}
\tightlist
\item
  Let A, B be two k-algebras (associative and unital), M a left A-module, P a left B-module, N a (B, A)-bimodule.
\end{enumerate}

\begin{enumerate}
\def\labelenumi{\alph{enumi})}
\tightlist
\item
  Show that there exist two spectral sequences (X, pp.~175-177, exercises 14 to 17) of \(k\) -modules \({}^{1}\mathrm{E}\) and \({}^{2}\mathrm{E}\) , converging to the same \(k\) -graded module, such that
\end{enumerate}

\[
{ } ^ { 1 } \mathrm{E} _ { 2 } ^ { p q } = \operatorname{Ext} _ { \mathbf { A } } ^ { p } ( \mathbf { M } , \operatorname{Ext} _ { \mathbf { B } } ^ { q } ( \mathbf { N } , \mathbf { P } ) ) ; \quad { } ^ { 2 } \mathrm{E} _ { 2 } ^ { p q } = \operatorname{Ext} _ { \mathbf { B } } ^ { p } ( \operatorname{Tor} _ { q } ^ { \mathbf { A } } ( \mathbf { N } , \mathbf { M } ) , \mathbf { P } ) .
\]

\begin{enumerate}
\def\labelenumi{\alph{enumi})}
\setcounter{enumi}{1}
\tightlist
\item
  By composing the homomorphisms \(\gamma\) and \(\delta\) of Exercise 15, p.~176, define a \(k\) -homomorphism
\end{enumerate}

\[
\mu : \operatorname{Ext} _ {\mathbf {A}} (\mathbf {M}, \operatorname{Hom} _ {\mathbf {B}} (\mathbf {N}, \mathbf {P})) \rightarrow \operatorname{Hom} _ {\mathbf {B}} \left(\operatorname{Tor} ^ {\mathbf {A}} (\mathbf {N}, \mathbf {M}), \mathbf {P}\right).
\]

If P is injective, prove that μ is an isomorphism.

\begin{enumerate}
\def\labelenumi{\alph{enumi})}
\setcounter{enumi}{2}
\item
  If N is a flat A-module, show that the sequence \({}^{1}E\) converges to Ext \(_{B}(N \otimes_{A} M, P)\) ; if N is a projective B-module, the sequence \({}^{2}E\) converges to Ext \(_{A}(M, Hom_{B}(N, P))\) .
\item
  Let \(\rho: A \to B\) a homomorphism of \(k\) -algebras. Show that there exists a spectral sequence E, such that \(\mathbf{E}_2^{pq} = \operatorname{Ext}_{\mathbf{B}}^p (\operatorname{Tor}_q^A(B, M), P)\) , converging to \(\operatorname{Ext}_{\mathbf{A}}(M, P)\) and a spectral sequence 'E, such that
\end{enumerate}

\[
{ } ^ { \prime } \mathrm{E} _ { 2 } ^ { p q } = \operatorname{Ext} _ { \mathbf { B } } ^ { p } ( \mathbf { P } , \operatorname{Ext} _ { \mathbf { A } } ^ { q } ( \mathbf { B } , \mathbf { M } ) ) ,
\]

converging to \(\operatorname{Ext}_{\mathbf{A}}(\mathbf{P}, \mathbf{M})\) .

\begin{enumerate}
\def\labelenumi{\arabic{enumi})}
\setcounter{enumi}{6}
\tightlist
\item
  Let C be an A-complex and M an A-module. One calls the module of extensions of M by C (resp. of C by M) the graded k-module
\end{enumerate}

\[
\mathcal {E} x t _ {\mathrm{A}} (\mathrm{C}, \mathrm{M}) = \mathrm{H} (\operatorname{Homgr} _ {\mathrm{A}} (\mathrm{C}, \mathrm{I} (\mathrm{M}))) \quad \left(\text { resp. } \mathcal {E} x t _ {\mathrm{A}} (\mathrm{M}, \mathrm{C}) = \mathrm{H} (\operatorname{Homgr} _ {\mathrm{A}} (\mathrm{L} (\mathrm{M}), \mathrm{C}))\right).
\]

If \(f: \mathbf{C} \to \mathbf{C}'\) is a morphism of complexes and \(g: \mathbf{M} \to \mathbf{M}'\) a homomorphism of A-modules, we set \(\mathcal{E}xt_{\mathbf{A}}(f, g) = \mathrm{H}(\mathrm{Homgr}(f, \mathrm{l}(g)))\) , \(\mathcal{E}xt_{\mathbf{A}}(g, f) = \mathrm{H}(\mathrm{Homgr}(\mathrm{L}(g), f))\) .

\begin{enumerate}
\def\labelenumi{\alph{enumi})}
\item
  If \(f: \mathbf{C} \to \mathbf{C}'\) is a homologism, the \(k\) -homomorphisms \(\mathcal{E}xt_{\mathbf{A}}(f, 1)\) and \(\mathcal{E}xt_{\mathbf{A}}(1, f)\) are bijective. If C is projective and bounded on the right (resp. injective and bounded on the left), the \(k\) -graded module \(\mathcal{E}xt_{\mathbf{A}}(\mathbf{C}, \mathbf{M})\) (resp. \(\mathcal{E}xt_{\mathbf{A}}(\mathbf{M}, \mathbf{C})\) ) is isomorphic to \(\mathrm{H}(\mathrm{Homgr}_{\mathbf{A}}(\mathbf{C}, \mathbf{M}))\) (resp. \(\mathrm{H}(\mathrm{Homgr}_{\mathbf{A}}(\mathbf{M}, \mathbf{C}))\)) .
\item
  Let \(0 \to C' \to C \to C'' \to 0\) be an exact sequence of A-complexes and \(0 \to M' \to M \to M'' \to 0\) an exact sequence of A-modules. Show that there exist exact sequences
\end{enumerate}

\[
\begin{array}{l} \dots \to \mathcal {E} x t _ {\mathrm{A}} ^ {n} (\mathrm{C} ^ {\prime \prime}, \mathrm{M}) \to \mathcal {E} x t _ {\mathrm{A}} ^ {n} (\mathrm{C}, \mathrm{M}) \to \mathcal {E} x t _ {\mathrm{A}} ^ {n} (\mathrm{C} ^ {\prime}, \mathrm{M}) \to \mathcal {E} x t _ {\mathrm{A}} ^ {n + 1} (\mathrm{C} ^ {\prime \prime}, \mathrm{M}) \to \dots \\ \dots \to \mathcal {E} x t _ {\mathrm{A}} ^ {n} (\mathrm{C}, \mathrm{M} ^ {\prime}) \to \mathcal {E} x t _ {\mathrm{A}} ^ {n} (\mathrm{C}, \mathrm{M}) \to \mathcal {E} x t _ {\mathrm{A}} ^ {n} (\mathrm{C}, \mathrm{M} ^ {\prime \prime}) \to \mathcal {E} x t _ {\mathrm{A}} ^ {n + 1} (\mathrm{C}, \mathrm{M} ^ {\prime}) \to \dots \\ \dots \to \mathcal {E} x t _ {\mathrm{A}} ^ {n} (\mathrm{M}, \mathrm{C} ^ {\prime}) \to \mathcal {E} x t _ {\mathrm{A}} ^ {n} (\mathrm{M}, \mathrm{C}) \to \mathcal {E} x t _ {\mathrm{A}} ^ {n} (\mathrm{M}, \mathrm{C} ^ {\prime \prime}) \to \mathcal {E} x t _ {\mathrm{A}} ^ {n + 1} (\mathrm{M}, \mathrm{C} ^ {\prime}) \to \dots \\ \dots \to \mathcal {E} x t _ {\mathrm{A}} ^ {n} (\mathrm{M} ^ {\prime \prime}, \mathrm{C}) \to \mathcal {E} x t _ {\mathrm{A}} ^ {n} (\mathrm{M}, \mathrm{C}) \to \mathcal {E} x t _ {\mathrm{A}} ^ {n} (\mathrm{M} ^ {\prime}, \mathrm{C}) \to \mathcal {E} x t _ {\mathrm{A}} ^ {n + 1} (\mathrm{M} ^ {\prime \prime}, \mathrm{C}) \to \dots \end{array}
\]

\begin{enumerate}
\def\labelenumi{\alph{enumi})}
\setcounter{enumi}{2}
\tightlist
\item
  Show that there exists a spectral sequence ``E (X, pp.~175-177, exercises 14 to 17) converging to \(\mathcal{E}xt_{\mathrm{A}}(\mathrm{C},\mathrm{M})\) such that \(E_2^{pq} = H^q (\mathrm{Ext}_A^p (\mathrm{C},\mathrm{M}))\) . If C is bounded above, show that there exists a spectral sequence 'E, converging to \(\mathcal{E}xt_{\mathrm{A}}(\mathrm{C},\mathrm{M})\) , such that \(E_2^{pq} = \mathrm{Ext}_A^p (\mathrm{H}_q(\mathrm{C}),\mathrm{M})\) (consider the bicomplex \(\mathrm{Homgr}_{\mathrm{A}}(\mathrm{C},\mathrm{I}(\mathrm{M}))\) ).
\end{enumerate}

\begin{enumerate}
\def\labelenumi{\arabic{enumi})}
\setcounter{enumi}{7}
\tightlist
\item
  Let P, Q be two A-complexes. Suppose either that P is bounded below and Q bounded above, or that there exists an integer d such that every A-module admits a projective resolution of length \(\leqslant d\) .
\end{enumerate}

\begin{enumerate}
\def\labelenumi{\alph{enumi})}
\tightlist
\item
  Show that there exist two spectral sequences 'E and ``E, converging to the same graded k-module, such that
\end{enumerate}

\[
^ {\prime} \mathrm{E} _ {2} ^ {p q} = \mathrm{H} ^ {p} (\mathrm{Ext} _ {\mathrm{A}} ^ {q} (\mathrm{P}, \mathrm{Q})) \quad^ {\prime \prime} \mathrm{E} _ {2} ^ {p q} = \bigoplus_ {q ^ {\prime} + q ^ {\prime \prime} = q} \mathrm{Ext} _ {\mathrm{A}} ^ {p} (\mathrm{H} _ {q ^ {\prime}} (\mathrm{P}), \mathrm{H} ^ {q ^ {\prime \prime}} (\mathrm{Q}))
\]

\((\text{ou} \operatorname{Ext}_{\mathbf{A}}^{q}(\mathbf{P}, \mathbf{Q}) \text{ denotes the complex associated with the bicomplex } \oplus \operatorname{Ext}_{\mathbf{A}}^{q}(\mathbf{P}_{r}, \mathbf{Q}^{s}))\) .

\begin{enumerate}
\def\labelenumi{\alph{enumi})}
\setcounter{enumi}{1}
\tightlist
\item
  If P is projective or Q is injective, the spectral sequence ``E converges to H(Homgr \(_A\) (P, Q)). If H(P) is projective or H(Q) is injective, the spectral sequence 'E converges to Homgr \(_A\) (H(P), H(Q)).
\end{enumerate}

